% -----------------------------------------------------------------------------------------------------
\section{一圈形状因子}\label{sec:sudakov-oneloop}

一圈结果显示了双极点 $1/\varepsilon^2$ 的来源（既软又共线的胶子），给出了后面全部全阶结构在最低阶的具体内容，
并通过向类时区的解析延拓暴露出 $\pi^2$ 增强。

\begin{definition}[夸克形状因子]\label{def:sudakov-ff}
令 $j^\mu=\bar\psi\gamma^\mu\psi$ 为单一无质量夸克味的矢量流（单位电荷）。对动量为 $p_1$ 的入射在壳夸克与动量为 $p_2$ 的出射在壳夸克，
$p_1^2=p_2^2=0$，定义顶点函数 $\Gamma^\mu$ 与\emph{类空形状因子} $F$：
\begin{equation}\label{eq:sudakov-ff-def}
\begin{gathered}
\langle q(p_2)|\,j^\mu(0)\,|q(p_1)\rangle=\bar u(p_2)\,\Gamma^\mu(p_1,p_2)\,u(p_1),\qquad
\Gamma^\mu=F(q^2)\,\gamma^\mu,\\
q=p_2-p_1,\qquad Q^2\equiv-q^2=2p_1\cdot p_2>0 ;
\end{gathered}
\end{equation}
色与自旋指标已从两边约去（$F$ 是色单位矩阵的系数）。对产生过程 $\gamma^*(q)\to q(p_1)\,\bar q(p_2)$，
$\langle q(p_1)\bar q(p_2)|j^\mu(0)|0\rangle=\bar u(p_1)\,F(q^2)\gamma^\mu\,v(p_2)$，$q=p_1+p_2$，$q^2=s>0$；
\emph{类时形状因子}是同一个解析函数在 $q^2=s+\ii0$ 的边界值。$F$ 在微扰论中按 $\alpha_s=\alpha_s(\mu)$ 展开，$F=1+F^{(1)}+O(\alpha_s^2)$。
\uses{def:prelim-feynman-rules}
\end{definition}

只出现一个形状因子的理由：无质量 QCD 的 Feynman 规则只产生含奇数个 $\gamma$ 矩阵的 Dirac 串，夹在 $\bar u(p_2)\cdots u(p_1)$ 之间时，
任何 $\sla p_1$、$\sla p_2$ 都可用反对易关系移到相应旋量旁而消失，只剩 $\gamma^\mu$ 乘以标量函数；这只用到 $\{\gamma^\mu,\gamma^\nu\}=2g^{\mu\nu}$，在 CDR 中逐阶成立。

\begin{proposition}[一圈形状因子（精确到 $\varepsilon$ 的所有阶） \tagR]\label{prop:sudakov-oneloop}
（$R_{\mathrm{PT}}$：在重整化微扰论中对红外正规化的部分子态严格。）在 $O(\alpha_s)$、$\MSbar$ 方案、$D=4-2\varepsilon$ 维中，类空形状因子为
\begin{equation}\label{eq:sudakov-F1}
\begin{gathered}
F(Q^2)=1+\frac{\alpha_s}{4\pi}C_F\Big(\frac{\mu^2}{Q^2}\Big)^{\varepsilon}f(\varepsilon)+O(\alpha_s^2),\\
f(\varepsilon)=-K_v(\varepsilon)\,\frac{2-\varepsilon+2\varepsilon^2}{\varepsilon^2(1-2\varepsilon)},\qquad
K_v(\varepsilon)=\frac{\ee^{\varepsilon\gamma_E}\Gamma(1+\varepsilon)\Gamma^2(1-\varepsilon)}{\Gamma(1-2\varepsilon)},
\end{gathered}
\end{equation}
对 $\varepsilon$ 精确；其 Laurent 展开为
\begin{equation}\label{eq:sudakov-F1-exp}
f(\varepsilon)=-\frac2{\varepsilon^2}-\frac3\varepsilon-8+\zeta_2+\varepsilon\Big(-16+\frac32\zeta_2+\frac{14}3\zeta_3\Big)+O(\varepsilon^2)
% VF-SUD-1
\end{equation}
（$\zeta_2=\pi^2/6$）。(a) 式中所有极点都是红外极点。(b) 类时区的值由 $(\mu^2/Q^2)^\varepsilon\to(\mu^2/s)^\varepsilon\ee^{\ii\pi\varepsilon}$ 得到：
\begin{equation}\label{eq:sudakov-F1-timelike}
\begin{gathered}
2\Real F^{(1)}(s+\ii0)=\frac{\alpha_sC_F}{2\pi}\Big(\frac{\mu^2}{s}\Big)^{\varepsilon}\cos(\pi\varepsilon)\,f(\varepsilon),\\
2\Real F^{(1)}(s+\ii0)-2F^{(1)}(Q^2\!=\!s)=\frac{\alpha_sC_F}{2\pi}\Big(\frac{\mu^2}{s}\Big)^{\varepsilon}\Big[\pi^2+\frac{3\pi^2}{2}\varepsilon+O(\varepsilon^2)\Big],
% VF-SUD-2
\end{gathered}
\end{equation}
即在 $|q^2|$ 相同时，类时值比类空值大 $\frac{\alpha_sC_F}{2\pi}\pi^2$，且这一差是红外有限的。（本命题只涉及部分子态。）
\uses{def:sudakov-ff,def:qcd-dimreg,def:prelim-dr-conventions,lem:prelim-scaleless,ex:prelim-scaleless,def:prelim-feynman-rules,def:prelim-msbar}
\end{proposition}

\begin{proof}
\emph{第一步：图与 Feynman 规则。}在 $O(\alpha_s)$，只有一个单粒子不可约图：入射夸克发出胶子 $k$、吸收光子、再吸收该胶子（\cref{fig:sudakov-ff}左图画的是同一个图的类时形式）。
外线自能插入是无质量在壳的、没有标度的积分，在维数正规化中为零（\cref{lem:prelim-scaleless}；其紫外/红外记账见\cref{ex:prelim-scaleless}，下面第五步再用）。
按\cref{def:prelim-feynman-rules}（Feynman 规范），
\begin{equation}\label{eq:sudakov-vertex}
\bar u(p_2)\Lambda^\mu u(p_1),\qquad
\Lambda^\mu=\int\!\frac{\dd^Dk}{(2\pi)^D}\,(\ii g\gamma^\nu t^a)\,\frac{\ii(\sla p_2-\sla k)}{(p_2-k)^2}\,\gamma^\mu\,\frac{\ii(\sla p_1-\sla k)}{(p_1-k)^2}\,(\ii g\gamma_\nu t^a)\,\frac{-\ii}{k^2},
\end{equation}
各分母带 $+\ii0$。常数因子为 $(\ii g)^2\cdot\ii\cdot\ii\cdot(-\ii)=-\ii g^2$，色因子 $t^at^a=C_F$，故
\[
\Lambda^\mu=-\ii g^2C_F\int\!\frac{\dd^Dk}{(2\pi)^D}\,\frac{N^\mu}{D_0D_1D_2},\qquad N^\mu=\gamma^\nu(\sla p_2-\sla k)\gamma^\mu(\sla p_1-\sla k)\gamma_\nu ,
\]
$D_0=k^2$，$D_1=(p_1-k)^2$，$D_2=(p_2-k)^2$。

\emph{第二步：Feynman 参数与分子代数。}用 $\frac1{D_0D_1D_2}=2\int\!\dd F\,(zD_0+xD_1+yD_2)^{-3}$，$\int\!\dd F\equiv\int_0^1\dd x\,\dd y\,\dd z\,\delta(1-x-y-z)$。
由 $p_i^2=0$ 与 $2p_1\cdot p_2=Q^2$，$zD_0+xD_1+yD_2=(k-xp_1-yp_2)^2-\Delta$，$\Delta=xyQ^2$。令 $\ell=k-xp_1-yp_2$，则
$\sla p_2-\sla k=\sla a-\sla\ell$、$\sla p_1-\sla k=\sla b-\sla\ell$，$a=(1-y)p_2-xp_1$，$b=(1-x)p_1-yp_2$。$\ell$ 的奇次项积分为零，$\ell^\alpha\ell^\beta\to g^{\alpha\beta}\ell^2/D$。
用 $D$ 维恒等式 $\gamma^\nu\sla a\gamma^\mu\sla b\gamma_\nu=-2\sla b\gamma^\mu\sla a+(4-D)\sla a\gamma^\mu\sla b$ 与 $\gamma^\rho\gamma^\mu\gamma_\rho=-(D-2)\gamma^\mu$，
\[
N^\mu\ \to\ -2\sla b\gamma^\mu\sla a+(4-D)\sla a\gamma^\mu\sla b+\frac{(D-2)^2}{D}\,\ell^2\gamma^\mu .
\]
夹在旋量之间：$\bar u(p_2)\sla b=(1-x)\bar u(p_2)\sla p_1$，$\sla a u(p_1)=(1-y)\sla p_2u(p_1)$，而
$\bar u(p_2)\sla p_1\gamma^\mu\sla p_2u(p_1)=-2p_1\cdot p_2\,\bar u(p_2)\gamma^\mu u(p_1)$（把 $\sla p_1$ 右移、$\sla p_2$ 左移，用 Dirac 方程）。于是
$\bar u\sla b\gamma^\mu\sla a u=-(1-x)(1-y)Q^2\,\bar u\gamma^\mu u$，同理 $\bar u\sla a\gamma^\mu\sla b u=-xyQ^2\,\bar u\gamma^\mu u$，
\[
\bar u(p_2)N^\mu u(p_1)\ \to\ \Big\{\big[2(1-x)(1-y)-2\varepsilon xy\big]Q^2+\frac{4(1-\varepsilon)^2}{D}\,\ell^2\Big\}\,\bar u(p_2)\gamma^\mu u(p_1).
\]

\emph{第三步：圈积分。}标准公式（\cref{def:qcd-dimreg}；附录~C 的\cref{sec:appC-formfactor}）
\[
\int\!\frac{\dd^D\ell}{(2\pi)^D}\frac{1}{(\ell^2-\Delta)^3}=\frac{-\ii}{(4\pi)^{D/2}}\frac{\Gamma(1+\varepsilon)}{2}\Delta^{-1-\varepsilon},\qquad
\int\!\frac{\dd^D\ell}{(2\pi)^D}\frac{\ell^2}{(\ell^2-\Delta)^3}=\frac{\ii}{(4\pi)^{D/2}}\frac{D}{4}\Gamma(\varepsilon)\Delta^{-\varepsilon}
\]
给出（$-\ii\cdot2\cdot(-\ii/2)=-1$，$-\ii\cdot2\cdot\ii/4=1/2$）
\[
F^{(1)}=\frac{2g^2C_F}{(4\pi)^{2-\varepsilon}}\,\Gamma(1+\varepsilon)\,(Q^2)^{-\varepsilon}\int\!\dd F\,
\Big\{\underbrace{\frac{(1-\varepsilon)^2}{\varepsilon}(xy)^{-\varepsilon}}_{\text{紫外}}-\underbrace{\big[(1-x)(1-y)-\varepsilon xy\big](xy)^{-1-\varepsilon}}_{\text{红外}}\Big\}.
\]

\emph{第四步：参数积分。}用 Dirichlet 积分 $\int\!\dd F\,x^{\alpha-1}y^{\beta-1}z^{\gamma-1}=\Gamma(\alpha)\Gamma(\beta)\Gamma(\gamma)/\Gamma(\alpha+\beta+\gamma)$
与 $(1-x)(1-y)=z+xy$（因 $1-x=y+z$，$1-y=x+z$），花括号的积分为
\[
\begin{aligned}
&\frac{(1-\varepsilon)^2}{\varepsilon}\frac{\Gamma^2(1-\varepsilon)}{\Gamma(3-2\varepsilon)}-\frac{\Gamma^2(-\varepsilon)}{\Gamma(2-2\varepsilon)}-(1-\varepsilon)\frac{\Gamma^2(1-\varepsilon)}{\Gamma(3-2\varepsilon)}\\
&\qquad=\frac{\Gamma^2(1-\varepsilon)}{\Gamma(1-2\varepsilon)}\,\frac{1}{1-2\varepsilon}\Big[\frac{1-\varepsilon}{2\varepsilon}-\frac{1}{\varepsilon^2}-\frac12\Big]
=-\frac{\Gamma^2(1-\varepsilon)}{\Gamma(1-2\varepsilon)}\,\frac{2-\varepsilon+2\varepsilon^2}{2\varepsilon^2(1-2\varepsilon)},
\end{aligned}
\]
其中用了 $\Gamma(-\varepsilon)=-\Gamma(1-\varepsilon)/\varepsilon$、$\Gamma(3-2\varepsilon)=2(1-\varepsilon)(1-2\varepsilon)\Gamma(1-2\varepsilon)$。最后代入
$\MSbar$ 的裸耦合 $g_0^2=4\pi\alpha_s\mu^{2\varepsilon}\ee^{\varepsilon\gamma_E}(4\pi)^{-\varepsilon}+O(\alpha_s^2)$（\cref{def:prelim-msbar}），
$g_0^2/(4\pi)^{2-\varepsilon}=\frac{\alpha_s}{4\pi}\mu^{2\varepsilon}\ee^{\varepsilon\gamma_E}$，即得式~\eqref{eq:sudakov-F1}。式~\eqref{eq:sudakov-F1-exp} 由
$\ln K_v(\varepsilon)=-\frac{\zeta_2}2\varepsilon^2-\frac73\zeta_3\varepsilon^3+O(\varepsilon^4)$（由 $\ln\Gamma(1+z)=-\gamma_Ez+\sum_{k\ge2}(-1)^k\zeta_kz^k/k$）与
$\frac{2-\varepsilon+2\varepsilon^2}{1-2\varepsilon}=2+3\varepsilon+8\varepsilon^2+16\varepsilon^3+O(\varepsilon^4)$ 相乘得到。

\emph{第五步：极点的来源（(a)）。}第三步中标为“紫外”的项来自 $\ell^2\to\infty$，它的极点 $+\frac{\alpha_sC_F}{4\pi}\frac1{\varepsilon_{\rm UV}}$ 是顶点的紫外发散；
由 Ward 恒等式 $Z_1=Z_2$，它与夸克场重整化的紫外部分相消。外线自能 $\delta Z_2^{\rm OS}=\frac{\alpha_sC_F}{4\pi}\big(-\frac1{\varepsilon_{\rm UV}}+\frac1{\varepsilon_{\rm IR}}\big)=0$
在维数正规化中为零（\cref{ex:prelim-scaleless}），把这一紫外极点“转换”为红外极点（\cref{def:prelim-dr-conventions}(a)）。守恒流不需要重整化，所以 $F$ 中余下的全部极点是红外极点：
双极点来自既软又共线的胶子，$-3/\varepsilon$ 中含共线极点与上述转换来的部分。

\emph{第六步：类时区（(b)）。}对 $\gamma^*\to q(p_1)\bar q(p_2)$ 重复第一至第四步（旋量换成 $\bar u(p_1)\cdots v(p_2)$），唯一的变化是 $\Delta=-xy(q^2+\ii0)$，
$+\ii0$ 来自 Feynman 传播子。对 $q^2=s>0$，$-s-\ii0=s\,\ee^{-\ii\pi}$，故 $(xyQ^2)^{-\varepsilon}\to(xys)^{-\varepsilon}\ee^{\ii\pi\varepsilon}$，从而
$F^{(1)}(s+\ii0)=\frac{\alpha_sC_F}{4\pi}(\mu^2/s)^\varepsilon\ee^{\ii\pi\varepsilon}f(\varepsilon)$。取实部得式~\eqref{eq:sudakov-F1-timelike} 的第一式；
差值为 $\frac{\alpha_sC_F}{2\pi}(\mu^2/s)^\varepsilon(\cos\pi\varepsilon-1)f(\varepsilon)$，由 $\cos\pi\varepsilon-1=-\frac{\pi^2\varepsilon^2}2+O(\varepsilon^4)$ 与 $f=-\frac2{\varepsilon^2}-\frac3\varepsilon+\cdots$
得方括号 $\pi^2+\frac32\pi^2\varepsilon+O(\varepsilon^2)$：$f$ 的 $\varepsilon^{-2}$ 与 $\varepsilon^{-1}$ 项都被 $\varepsilon^2$ 乘掉，差值有限。
\end{proof}

把式~\eqref{eq:sudakov-F1-exp} 乘上 $(\mu^2/Q^2)^\varepsilon$ 并展开（$L\equiv\ln(Q^2/\mu^2)$），
\begin{equation}\label{eq:sudakov-F1-L}
F^{(1)}(Q^2)=\frac{\alpha_sC_F}{4\pi}\Big[-\frac2{\varepsilon^2}+\frac{2L-3}{\varepsilon}-L^2+3L-8+\zeta_2\Big]+O(\varepsilon);
% VF-EE-4
\end{equation}
类时区以 $L\to\ell-\ii\pi$（$\ell\equiv\ln(s/\mu^2)$）代入得 $2\Real F^{(1)}(s+\ii0)=\frac{\alpha_sC_F}{2\pi}\big[-\frac2{\varepsilon^2}+\frac{2\ell-3}{\varepsilon}-\ell^2+3\ell-8+\zeta_2+\pi^2\big]$，虚部见\cref{prop:sudakov-coulomb}。
$f(\varepsilon)$ 直到 $\varepsilon^4$ 与\cite[式~(3.5)]{MochVermaserenVogt2005}的单圈系数逐项相同（两圈见\cite{Gonsalves1983}；上述每一步都经符号与数值计算核对）。% check: s1_formfactor.py
取 $\alpha_s=0.1180$（世界平均值\cite{ParticleDataGroup:2024cfk}，仅作示意），$\frac{\alpha_sC_F}{2\pi}\pi^2\approx0.247$：解析延拓本身给出约 $25\%$ 的修正，
它在 Drell--Yan 系数函数的 $\delta(1-z)$ 项中重现（第~12~章\cref{prop:dy-nlo}），其指数化见\cref{cor:sudakov-exponentiation}。
